% Near-verbatim source section for OY_napMPywE, 01:05:00--01:15:00.
% YIN-VERBATIM-SOURCE-SHA256: 5ac8ac5fb25a3235d8fa11b2b6be99b5f2bb9329d307c4045629544f4e43e9bd
% Equation text remains subordinate to the cited NOTES_EXACT placeholders.

% YIN-VERBATIM: id=YIN-OY-T000239; time=01:05:00.000--01:05:14.579; speaker=Yin; notes=253a:6; pdf=11; disposition=unresolved_audio; uncertainty=Words around "if it's one thought" and the subject of "has to be independent" are unclear.
Yeah. This, if it's \textit{[unresolved]}, something infinitesimal,
\textit{[unresolved]} has to be independent.

% YIN-VERBATIM: id=YIN-OY-T000240; time=01:05:14.579--01:05:37.380; speaker=Yin; notes=253a:6; pdf=11; disposition=included_with_unresolved_audio; uncertainty=The words after "recover the finite form by" are unclear.
So I'm considering the case that these are small, and so we'll expand to first
order. This will ultimately be sufficient for our discussion, because I can
recover the finite form by \textit{[unresolved]}. All right.

% YIN-VERBATIM: id=YIN-OY-T000241; time=01:05:37.380--01:05:43.200; speaker=Yin; notes=253a:6; pdf=11; disposition=included_question_management; uncertainty=none.
Okay, so I can only do this. Any questions so far?

% YIN-VERBATIM: id=YIN-OY-T000242; time=01:05:43.200--01:06:12.200; speaker=mixed; notes=253a:6; pdf=11; disposition=student_question_unresolved; uncertainty=The question beginning "what you wrote here is just the general form for" has no intelligible ending.
Yes. \textit{[Question unresolved.]} Okay. Yeah.

% YIN-NOTE-EQUATION-PLACEHOLDER: key=infinitesimal-poincare-unitary; notes=253a:6; pdf=11; authority=NOTES_EXACT; placement=before generator names.
% YIN-VERBATIM: id=YIN-OY-T000243; time=01:06:12.200--01:06:24.359; speaker=Yin; notes=253a:6; pdf=11; disposition=included; uncertainty=none.
They have some names. So this $\hat P_\mu$ is the so-called energy-momentum
operator.

% YIN-VERBATIM: id=YIN-OY-T000244; time=01:06:24.359--01:06:38.400; speaker=Yin; notes=253a:6; pdf=11; disposition=included; uncertainty=none.
And the $\hat J^{\mu\nu}$ is the operator corresponding to the Lorentz boosts
and the angular momentum.

% YIN-VERBATIM: id=YIN-OY-T000245; time=01:06:38.400--01:07:00.960; speaker=Yin; notes=253a:6; pdf=11; disposition=unresolved_audio; uncertainty=The example following "if I" is missing.
So indeed, the claim is going to be that, for example, if I
\textit{[unresolved]}.

% YIN-VERBATIM: id=YIN-OY-T000246; time=01:07:00.960--01:07:33.240; speaker=Yin; notes=253a:6; pdf=11; disposition=unresolved_audio; uncertainty=The indexed quantity after "indices of" is unclear.
So you have to be careful with raising and lowering indices of
\textit{[unresolved]}.

% YIN-VERBATIM: id=YIN-OY-T000247; time=01:07:33.240--01:07:45.359; speaker=Yin; notes=253a:6; pdf=11; disposition=included_continues_next; uncertainty=none.
Now I'm not going to derive
% YIN-VERBATIM: id=YIN-OY-T000248; time=01:07:45.359--01:08:04.619; speaker=Yin; notes=253a:6; pdf=11; disposition=included; uncertainty=none.
the commutators of these $P$'s and $J$'s, using the property that the $U$'s
compose according to the rule of Poincare transformations.

% YIN-VERBATIM: id=YIN-OY-T000249; time=01:08:04.619--01:08:18.960; speaker=Yin; notes=253a:6; pdf=11; disposition=included_continues_next; uncertainty=none.
In particular, $U(\Lambda',a')$, let's say, composed with $U(\Lambda,a)$.
Well, this sends $x$ to $\Lambda x+a$, and
% YIN-VERBATIM: id=YIN-OY-T000250; time=01:08:18.960--01:08:28.920; speaker=Yin; notes=253a:6; pdf=11; disposition=included; uncertainty=none.
this one does it again with $\Lambda'$ and $a'$. This gives
$U(\Lambda'\Lambda,\Lambda'a+a')$, multiplying the matrices.

% YIN-NOTE-EQUATION-PLACEHOLDER: key=poincare-group-composition; notes=253a:6; pdf=11; authority=NOTES_EXACT; placement=after spoken composition.
% YIN-VERBATIM: id=YIN-OY-T000251; time=01:08:28.920--01:08:47.880; speaker=Yin; notes=253a:6; pdf=11; disposition=included; uncertainty=none.
I will let you think for yourself how to derive from this the commutation
relations of the $P$'s and $J$'s. The point is that they are all dictated by
the requirement that these $U$'s have an interpretation as the unitary
operators that generate Poincare transformations.

% YIN-NOTE-EQUATION-PLACEHOLDER: key=poincare-algebra; notes=253a:6; pdf=11; authority=NOTES_EXACT; placement=after derivation comment.
% YIN-VERBATIM: id=YIN-OY-T000252; time=01:08:47.880--01:09:06.480; speaker=Yin; notes=253a:6; pdf=11; disposition=unresolved_audio; uncertainty=Only "correspondence is actually although" survives.
\textit{[Unresolved speech.]}

% YIN-VERBATIM: id=YIN-OY-T000253; time=01:09:06.480--01:09:18.779; speaker=Yin; notes=253a:6; pdf=11; disposition=included_with_unresolved_audio; uncertainty=The adjective before "statement" is unclear.
Well, this is the \textit{[unresolved]} statement that energy-momentum are the
conserved charges associated with symmetry in time and space.

% YIN-VERBATIM: id=YIN-OY-T000254; time=01:09:18.779--01:09:29.460; speaker=Yin; notes=253a:6; pdf=11; disposition=included; uncertainty=none.
We will give a more systematic discussion in the context of field theory,
using the Noether procedure, in a few lectures.

% YIN-VERBATIM: id=YIN-OY-T000255; time=01:09:29.460--01:09:47.279; speaker=Yin; notes=253a:6; pdf=11; disposition=classroom_logistics; uncertainty=The referenced author or book rendered "new member" is unclear; canonical cleaned_text is null.

% YIN-VERBATIM: id=YIN-OY-T000256; time=01:09:47.279--01:10:06.440; speaker=Yin; notes=253a:6; pdf=11; disposition=included; uncertainty=none.
The discussion of the procedure, which we'll come to a little bit later, will
allow you to express these $P$'s and the $J$'s in terms of spatial integrals
of the stress-energy tensor.

% YIN-VERBATIM: id=YIN-OY-T000257; time=01:10:06.440--01:10:22.500; speaker=Yin; notes=253a:6; pdf=11; disposition=included_with_unresolved_audio; uncertainty=The modifier before "stress-energy tensor" is unclear.
You understand exactly how these are related to the \textit{[unresolved]}
stress-energy tensor. You can read about that in Chapter 7.

% YIN-VERBATIM: id=YIN-OY-T000258; time=01:10:22.500--01:10:29.820; speaker=Yin; notes=253a:6; pdf=11; disposition=included_with_unresolved_audio; uncertainty=The class of systems after "any" is unclear.
We'll come to this later, but this is a general fact of any
\textit{[unresolved]}.

% YIN-VERBATIM: id=YIN-OY-T000259; time=01:10:29.820--01:10:52.830; speaker=mixed; notes=253a:6; pdf=11; disposition=student_question_unresolved; uncertainty=The question and the word rendered "stagger" are unclear.
Yes, so \textit{[student question unresolved]}. Well, if you have one plus
$i$ times some Hermitian operator, \textit{[unresolved]}.

% YIN-VERBATIM: id=YIN-OY-T000260; time=01:10:52.830--01:11:09.420; speaker=unresolved; notes=253a:6; pdf=11; disposition=unresolved_weak_room_audio; uncertainty=Possible continuation of the unitarity exchange; canonical cleaned_text is null.
% YIN-VERBATIM: id=YIN-OY-T000261; time=01:11:09.420--01:11:17.520; speaker=Yin; notes=253a:6; pdf=11; disposition=classroom_logistics; uncertainty=none; canonical cleaned_text is null.

% YIN-VERBATIM: id=YIN-OY-T000262; time=01:11:17.520--01:11:32.960; speaker=Yin; notes=253a:6; pdf=11; disposition=included_continues_next; uncertainty=none.
So the point of this little discussion is just to explain what it means for
$\hat\phi$ to be tied to the spacetime point $x$. Let me just reiterate: it's
this equation here that says the Poincare symmetry
% YIN-VERBATIM: id=YIN-OY-T000263; time=01:11:32.960--01:11:46.140; speaker=Yin; notes=253a:6; pdf=11; disposition=included; uncertainty=none.
transformation, which is realized by the unitary operator $U$, should act as
$U\hat\phi(x)U^{-1}=\hat\phi(\Lambda x+a)$, the field at the transformed
spacetime point.

% YIN-NOTE-EQUATION-PLACEHOLDER: key=scalar-field-poincare-covariance; notes=253a:6; pdf=11; authority=NOTES_EXACT; placement=after covariance recap.
% YIN-VERBATIM: id=YIN-OY-T000264; time=01:11:46.140--01:12:07.440; speaker=Yin; notes=253a:6,7; pdf=11,12; disposition=included; uncertainty=none.
Okay, that's that. But so far I haven't said anything about locality, so
that's going to be the key. A key property is going to be microcausality.

% YIN-VERBATIM: id=YIN-OY-T000265; time=01:12:07.440--01:12:25.380; speaker=Yin; notes=253a:7; pdf=12; disposition=included_continues_next; uncertainty=none.
In this particular example, we're going to demand
$[\hat\phi(x),\hat\phi(y)]=0$, where $x$ and $y$ are two spacetime points,
whenever
% YIN-VERBATIM: id=YIN-OY-T000266; time=01:12:25.380--01:12:38.480; speaker=Yin; notes=253a:7; pdf=12; disposition=included_continues_next; uncertainty=none.
whenever $x$ and $y$ are spacelike separated. That is,
% YIN-VERBATIM: id=YIN-OY-T000267; time=01:12:38.480--01:12:57.900; speaker=Yin; notes=253a:7; pdf=12; disposition=included_with_unresolved_audio; uncertainty=The phrase after "Lorentzian inner product" is unclear.
And you always want to use the convention $(x-y)^2$. This is the Lorentzian
inner product \textit{[unresolved]}. This is defined to be minus
$(x^0-y^0)^2$, the time difference, plus the spatial difference
$(\vec x-\vec y)^2$.

% YIN-NOTE-EQUATION-PLACEHOLDER: key=microcausality-and-mostly-plus-interval; notes=253a:7; pdf=12; authority=NOTES_EXACT; placement=after spoken definition.
% YIN-VERBATIM: id=YIN-OY-T000268; time=01:12:57.900--01:13:08.699; speaker=Yin; notes=253a:7; pdf=12; disposition=included; uncertainty=none.
So, i.e., this is positive.

% YIN-VERBATIM: id=YIN-OY-T000269; time=01:13:08.699--01:13:21.960; speaker=Yin; notes=253a:7; pdf=12; disposition=included; uncertainty=none.
Okay. Of course, as I said, I'm always working in units where the speed of
light is set to one. Okay.

% YIN-VERBATIM: id=YIN-OY-T000270; time=01:13:21.960--01:13:31.800; speaker=Yin; notes=253a:7; pdf=12; disposition=included; uncertainty=none.
So that's microcausality, at least applied to this local field operator
$\hat\phi$ we're talking about here.

% YIN-VERBATIM: id=YIN-OY-T000271; time=01:13:31.800--01:13:38.760; speaker=Yin; notes=253a:7; pdf=12; disposition=classroom_logistics; uncertainty=none; canonical cleaned_text is null.

% YIN-VERBATIM: id=YIN-OY-T000272; time=01:13:38.760--01:13:53.820; speaker=Yin; notes=253a:7; pdf=12; disposition=included_with_unresolved_audio; uncertainty=Two words before and within "constraint on the quantum system" are unclear.
No, it's not obvious at the moment, but we're back at the requirement of the
existence of an operator $\hat\phi$ tied to a spacetime point $x$ in that way.
\textit{[Unresolved]} it's an \textit{[unresolved]} constraint on the quantum
system. In fact, you can ask
% YIN-VERBATIM: id=YIN-OY-T000273; time=01:13:53.820--01:13:59.880; speaker=Yin; notes=253a:7; pdf=12; disposition=included_continues_next; uncertainty=none.
Okay, so
% YIN-VERBATIM: id=YIN-OY-T000274; time=01:13:59.880--01:14:19.440; speaker=Yin; notes=253a:7; pdf=12; disposition=included; uncertainty=none.
you can ask whether this applies to our system of free particles. Intuitively
it should, because we have non-interacting particles, and the particles
themselves propagate with the relativistic dispersion relation, so the system
must be causal.

% YIN-VERBATIM: id=YIN-OY-T000275; time=01:14:19.440--01:14:27.659; speaker=Yin; notes=253a:7; pdf=12; disposition=included; uncertainty=none.
But can you actually write down these local field operators in that theory?

% YIN-VERBATIM: id=YIN-OY-T000276; time=01:14:27.659--01:14:44.760; speaker=Yin; notes=253a:7; pdf=12; disposition=included; uncertainty=none.
Yeah. What about our model, the very simple quantum-mechanical model of free
relativistic particles?

% YIN-NOTE-EQUATION-PLACEHOLDER: key=free-relativistic-hamiltonian; notes=253a:7; pdf=12; authority=NOTES_EXACT; placement=with endpoint sentence.
% YIN-VERBATIM: id=YIN-OY-T000277; time=01:14:44.760--01:15:00.000; speaker=Yin; notes=253a:7; pdf=12; disposition=included_cut_at_endpoint; uncertainty=The lecture continues past 01:15:00 with the momentum operator; displayed creator uses the transcript notation layer.
So, for example, I have already written down the Hamiltonian, which is
$H=\int d^{D-1}\vec p\,\sqrt{\vec p^{\,2}+m^2}\,
a_{\vec p}^{\dagger}a_{\vec p}$.

% YIN-VERBATIM-COUNTS: eligible_cleaned_words=660; retained_cleaned_words=660; retained_ratio=1.000.
